\scanpage{38}
\begin{center}
INSTITUT\\
DES HAUTES ÉTUDES SCIENTIFIQUES\\
\rule{2.3em}{0.4pt}
\end{center}

\noindent\begin{minipage}{0.6\linewidth}
SGA 4 $\underline{\overline{\mathrm{XVIII}}}$\\[0.8em]
\hspace*{3em}\rlap{\begin{tikzpicture}[overlay,x=1em,y=1em]
\draw[line width=.35pt] (1.0,-.6) .. controls (-1.0,-.5) and (-.8,.7) .. (1.0,.9)
 .. controls (3.1,1.1) and (7.1,1.2) .. (7.8,2.7);
\end{tikzpicture}}\inserted{(Deligne)}
\end{minipage}

\medskip
\noindent Références pour SGA 5 $\underline{\overline{\mathrm{VII}}}$\quad
$\left\{\begin{array}{@{}l@{}}
\text{Dém de 3.\deleted{\unsure{6}}5.3.}\\[.4em]
\text{1.1\qquad b)}
\end{array}\right.$

\medskip
\noindent\begin{tikzpicture}[x=\linewidth,y=1em]
\draw[line width=.4pt] (0,0)--(1,.8);
\end{tikzpicture}

$\underline{\overline{\mathrm{V}}}$\qquad dire quelques mots sur relations\\
\hspace*{2.7em}avec $\underleftarrow{\mathrm{lim}}$ et ${\underleftarrow{\mathrm{lim}}}^{(p)}$ :

a) Sur un topos de \unsure{la} forme $\widehat C$,\\
\hspace*{1.4em}les foncteurs $H^n$ sont les $\mathop{\underleftarrow{\mathrm{lim}}}\nolimits_C^{(n)}$

b) Je \unsure{n'ai|n'y a} plus une telle interpré-\\
\hspace*{1.4em}tation pour un topos général,\\
\hspace*{1.4em}puisqu'on est obligé de prendre
\margin{\raisebox{-3.5em}[0pt][0pt]{\rule{.35pt}{5.5em}}\hspace{.35em}\begin{tabular}[t]{@{}l@{}}si ça\\marche \ldots\end{tabular}}
\[
 H^{n}(\deleted{\unsure{E|C}}\unsure{\widetilde C|\widetilde E},F)\simeq
 \mathop{\underrightarrow{\mathrm{lim}}}_{\mathcal U\in\mathcal R}
 \mathop{\underleftarrow{\mathrm{lim}}}_{\mathcal U}
 \deleted{\unsure{H^n|F^n}}\unsure{F^{n}/\mathcal U_{\varphi}|F/\mathcal U_{\varphi}|F/C_{/\varphi}}
\]

$\mathcal U$ crible variable

On a un recouvrement de l'objet final de $\widehat C$ par les objets de $C$

\bigskip
\noindent\hspace*{2.4em}$X\unsure{\times|+}Y$
% archive sheet number: 37
